\chapter{Introduction}
%The Standard Model (SM) of the particle physics describes the nature of the constituents of the matter and the fundamental forces 
The Standard Model (SM) of the particle physics gives successful explanation about the constituents of the physical matter and the fundamental forces 
through which they interact with each other and with composite particles. 
%The three fundamental forces (electromagnetic, strong and weak) are described by three dedicated
The electromagnetic, strong and weak forces are described by dedicated
theoretical descriptions Quantum Electrodynamics (QED), Quantum Chromodynamics (QCD) and  Quantum Flavourdynamics (QFD) which are 
encoded in the SM. 
Later on, electromagnetic and weak forces among these, were unified as electroweak (EWK) force at the order of 100~GeV unification energy scale. 
Sheldon Glashow, Abdus Salam, and Steven Weinberg had significant contribution in this milestone. 
The experimental verifications of the existence of the electroweak interaction were established by the 
In 1973, existence of neutral currents announced by the Gargamelle collaboration and then in 1981, the discovery of two massive gauge bosons ($\W$ and $\Z$) by UA1 and  UA2 collaborations verified the existence of EWK force. 
These 
boson should be massless (like photon in QED) according to Goldstone's theorem. To overcome this issue, spontaneous symmetry breaking was introduced 
that generates the masses of the particles by the Higgs mechanism and separates weak interaction from  electromagnetic (EM) 
%interaction \ref{Reference7,Reference8,Reference9,Reference10,Reference11,Reference12}. Higgs boson, a new scalar particle, was introduced by the theory being last 
interaction ~\cite{Reference7,Reference8,Reference9,Reference10,Reference11,Reference12}. Higgs boson, a new scalar particle, was introduced by the theory being last 
missing patch of the SM. This triggered the search for this boson of the SM at experimental facilities which lasted for several years. 
One of the main design and construction motivations of the Large Hadron Collider (LHC) at the CERN was also to look for this new particle
%within full mass range allowed 
given unitary and above the Large Electron-Positron (LEP) constraints, i.e from 114~GeV to 1~TeV. After a long search, the CMS and the ATLAS collaborations of LHC, announced the existence of the heavy scalar boson (the Higgs)
on 4 July 2012~\cite{Reference20,Reference21}. Immediately after the discovery, further tests to measure its properties were begun. Mass of the newly discovered particle, the only free parameter in Higgs sector, has been 
measured to be $m_H = 125.09\pm 0.21(stat.)\pm0.11(syst.)$~GeV~\cite{Reference22}. Other properties such as spin, parity~\cite{Reference21a} and couplings of this boson to different particles~\cite{Reference21b,Reference21c} have also been measured and observed to consistent with the predictions of SM. Therefore, the new particle very probable appears as the SM Higgs boson, affiliated with EWK spontaneous symmetry breaking. 
SM has  passed enormous experimental tests in last 50 years and no significant deviation has been seen yet. Experimental evidence of the Higgs bosonis the biggest achievement of the SM. Other success includes the experimental verification of its prior predictions about the $\W$ and $\Z$ gauge bosons, gluon, top and charm quarks. %The ratio of the masses of $\W$ and $\Z$ bosons was also found to be consistent with SM predictions. 
\par 
Although SM has passed numerous experimental tests with great precision, it is not considered to be theory of everything. There are many 
phenomena which are not explained by SM i.e Hierarchy problem, neutrino massess and existence of dark matter and dark energy etc. There are several theoretical extension posed in SM such as the Minimal Super Symmetrical Model(MSSM) and the Next-to-Minimal Super Symmetric Standard Model (NMSSM) or completely new 
theoretical descriptions, such as string theory and extra dimensions. More precise Higgs boson properties measurements are needed to capture any possible deviation
from SM predictions which can be a doorway to new physics. 
\par 
Model independent production cross section for a particular process is an important property of H boson, where the cross section gives the possibility of the occurrence of that process.  In this dissertation, measurements of the integrated or inclusive and differential fiducial cross sections for H boson production in $H \rightarrow 4\ell$ decay channel ($\ell = e, \mu$) in proton-proton collisions at $\sqrt{s}=13$~TeV are presented. 
These measurements provide a direct test of perturbative QCD  calculations in the Higgs sector of SM at the TeV energy scale. Any deviation of the integrated or the differential cross sections from SM expectations would  provide a direct hint to BSM contributions. Nearly similar measurements
have already been reported by the ATLAS experiment in  $H \rightarrow \gamma \gamma$~\cite{ATLAS:Hgg_8TeV,ATLAS:Hgg_8TeV_13TeV}, 
$H \rightarrow 4\ell$~\cite{ATLAS:H4l_8TeV,ATLAS:H4l_8TeV_13TeV} decay  channels, combination of both 
channels~\cite{ATLAS:H4l_Hgg_8TeV,ATLAS:H4l_Hgg_8TeV_13TeV} and by the CMS experiment in $H \rightarrow \gamma \gamma$ decay mode~\cite{CMS:Hgg_8TeV}. % also add reference for HIG-16-041 
\par
The Higgs cross section measurement helps to probe the structure, nature of the interaction of H boson with decaying particles.  Differential fiducial cross 
sections are measured using the full Run 2 CMS data recorded at 13~\TeV corresponding to several kinematic observables which are sensitive to the H boson production 
mechanisms: rapidity and transverse momentum of the H boson, associated jet multiplicity and transverse momentum of leading jet.  These differential distributions are sensitive to production mechanism of gluon fusion, colliding protons' PDFs, theoretical modeling of hard quarks radiations and relative contributions of different production mechanisms of H boson.
%The results are also produced for kinematic observables which are sensitive to the Higgs boson decay mechanisms: the invariant masses of the two lepton pairs and five helicity angles formed by the leptons in the Collins-Soper frame~\cite{Reference40}.

\par The fragmentation of the dissertation is  systematized as follows. An overview of the SMof particle physics, production and decays of H boson are described in Chapter ~\ref{Chapter1}. Overview of the LHC and its CMS experiment are presented in Chapter ~\ref{Chapter2}. The calibrations and reconstructions of the  physics objects within the CMS detector are
described in Chapter ~\ref{Chapter3}. 
Chapter ~\ref{Chapter4} presents the procedure of measurements of inclusive and differential fiducial cross sections of H boson production at $\sqrt{s} = 13$~TeV. 
%Chapter \ref{Chapter5} presents the measurements of differential and integrated fiducial cross sections for Z boson production at $\sqrt{s} = 8$~TeV. 
%Chapter \ref{Chapter6} presents the integrated measurement of integrated fiducial cross sections for Higgs boson production at $\sqrt{s} = 13$~TeV and Higgs boson production cross sections as a function of $\sqrt{s}$.
Finally, the conclusions and outlook of the dissertation are summarized in Chapter ~\ref{Chapter7}.
%\ref{Chapter7} and \ref{Chapter8}.
